\subsection{局部可积函数}

命题 1. --- 设 g 是 T 中局部几乎处处（关于正测度 \(\mu\) ）有定义的一个函数，取值于 Banach 空间 F（相应地，\(\overline{R}\)）中。下列性质等价：

\begin{enumerate}
\def\labelenumi{\alph{enumi})}
\item
  对每个点 \(t \in \mathbb{T}\) ，存在一个邻域 \(\mathbb{V}\) （\(t\) 的），使得函数 \(\mathbf{g}\varphi_{\mathbb{V}}\) 是 \(\mu\) -可积的。
\item
  函数 \(\mathbf{g}\) 是 \(\mu\) -可测的，且对每个紧集 \(\mathrm{K} \subset \mathrm{T}\) ，\(\int^{*} |\mathbf{g}| \varphi_{\mathrm{K}} d\mu < +\infty\) 。
\item
  对每个数值函数 \(h \in \mathcal{K}(\mathrm{T})\) ，gh 都是 \(\mu\) -可积的。
\end{enumerate}

让我们证明 a) 蕴涵 b)；因为，由局部化原理（Ch. IV, §5, No.~2, 命题 4），函数 g 是可测的。另一方面，对每个 \(t \in \mathbf{K}\) ，由假设存在一个邻域 \(V_t\) （\(t\) 的，在 \(T\) 中），使 \(\mathbf{g}\varphi_{V_t}\) 可积；因此 \(K\) 可以被有限个邻域 \(V_i\) （ \(1 \leqslant i \leqslant n\) ）覆盖，使诸函数 \(\mathbf{g}\varphi_{V_i}\) 都可积。由于 \(|\mathbf{g}| \varphi_K \leqslant \sum_{i=1}^{n} |\mathbf{g}| \varphi_{V_i}\) ，就有 \(\int^* |\mathbf{g}| \varphi_K d\mu < +\infty\) 。

其次，b) 蕴涵 c)，因为此时 gh 可测，且若 L 是 h 的紧支集，则 \(|gh| \leqslant \|h\|\cdot|g|\varphi_{L}\) ，故由假设 \(\int^{*}|gh|d\mu < +\infty\) ；于是由可积性判别法（Ch. IV, §5, No.~6, 定理 5），gh 可积。

最后，c) 蕴涵 a)。的确，对每个 \(t \in T\) ，取 V 为 t 的一个紧邻域。存在 T 到 {[}0,1{]} 中的一个连续映射 h，它在 V 上等于 1 且具有紧支集（Ch. III, §1, No.~2, 引理 1）；由假设 gh 可积，因此 \(\mathbf{g}\varphi_{\mathrm{V}} = (\mathbf{g}h)\varphi_{\mathrm{V}}\) 也可积（Ch. IV, §5, No.~6, 定理 5 的推论 3）。

定义 1. --- 在 T 中局部几乎处处（关于正测度 \(\mu\) ）有定义、取值于 Banach 空间 F（相应地，\(\overline{R}\)）中的函数 g，如果满足命题 1 的条件 a), b), c)，就称为关于 \(\mu\) 局部可积的（或局部 \(\mu\) -可积的）。如果 \(\theta\) 是一个复测度，局部 \(\theta\) -几乎处处有定义的函数 g 称为局部 \(\theta\) -可积的，如果它关于正测度 \(|\theta|\) 局部可积。

如果 \(\mathbf{g}\) 是局部 \(\theta\) -可积的，那么每个局部几乎处处等于 \(\mathbf{g}\) 的函数都是局部可积的。显然，两个局部可积函数之和是局部可积的。取值于 \(\mathbf{F}\) 、处处有定义且关于 \(\theta\) 局部可积的函数构成一个向量空间，记作 \(\mathcal{L}_{\mathrm{loc}}^{1}(\mathrm{T},\theta ;\mathrm{F})\) ；当 \(\mathbf{F} = \mathbf{R}\) 或 \(\mathbf{C}\) 时，如无歧义，常省去对 \(\mathbf{F}\) 的提及。这个空间将总是（除明确说明相反者外）赋以由半范数 \(\mathbf{g}\mapsto \int |\mathbf{g}\varphi_{\mathbf{K}}|d|\theta |\) 所定义的拓扑，其中 \(\mathbf{K}\) 跑遍 \(\mathbf{T}\) 的紧子集所成的集。相联的 Hausdorff 空间，即 \(\mathcal{L}_{\mathrm{loc}}^{1}(\mathrm{T},\theta ;\mathrm{F})\) 关于局部几乎处处为零的映射所成的子空间 \(\mathcal{N}_{\mathrm{F}}^{\infty}\) 的商，记作 \(\mathrm{L}_{\mathrm{loc}}^{1}(\mathrm{T},\theta ;\mathrm{F})\) 。空间 \(\mathrm{L}_{\mathrm{loc}}^{1}(\mathrm{T},\theta ;\mathrm{F})\) 与 \(\mathrm{L}_{\mathrm{loc}}^{1}(\mathrm{T},|\theta |; \mathrm{F})\) 是相同的。

可以证明，刚才定义的拓扑向量空间是完备的（习题 31）。

在每个紧集上本性有界的每个可测函数 g 都是局部可积的。对每个使 \(1 \leqslant p \leqslant +\infty\) 的数 p，每个函数 \(g \in L_{F}^{p}\) 都是局部可积的；的确，对每个函数 \(h \in \mathcal{K}(T)\) ，h 属于 \(L^{q}\) （其中 q 是与 p 共轭的指数），因此 gh 可积（Ch. IV, §6, No.~4, 定理 2 的推论 4）。

设 \(F, G, H\) 是三个 Banach 空间，\((\mathbf{u}, \mathbf{v}) \mapsto \Phi(\mathbf{u}, \mathbf{v})\) 是 \(F \times G\) 到 \(H\) 中的一个连续双线性映射。如果 \(f\) 局部可积且取值于 F，且 \(g \in L_{G}^{\infty}\) ，那么 \(\Phi(\mathbf{f}, \mathbf{g})\) 是局部可积的（Ch. IV, §6, No.~4, 定理 2 的推论 1）。

